\documentclass[journal]{IEEEtran}

\usepackage{amsmath,amsfonts,amssymb}
\usepackage{algorithmic}
\usepackage{graphicx}
\usepackage{textcomp}
\usepackage{xcolor}
\usepackage{booktabs}
\usepackage{cite}
\usepackage{url}
\usepackage{hyperref}

\hypersetup{
    colorlinks=true,
    linkcolor=blue,
    citecolor=blue,
    urlcolor=blue
}

\begin{document}

\title{A Decoupled Recurrent Actor--Critic Reinforcement Learning Framework for Congestion-Aware Smart Parking Allocation}

\author{Smart Parking BTP Research Team%
\thanks{Manuscript prepared November 2026 for academic publication. The authors are with the Department of Computer Science and Engineering.}}

\maketitle

\begin{abstract}
Urban cruising for parking generates severe traffic congestion, accounting for up to 30\% of localized vehicular delay in metropolitan cores. While static heuristic policies (e.g., First-In-First-Out, nearest-zone greedy allocation) remain ubiquitous, they fail under stochastic arrival bursts and uncertain parking dwell turnovers. In this work, we propose a novel 24,420-parameter Recurrent Actor--Critic architecture utilizing Proximal Policy Optimization (PPO) tailored for asynchronous minute-level parking guidance. Crucially, we formulate a decoupled arrival-conditioned learning mechanism: policy and entropy losses are conditioned strictly on discrete arrival decision epochs ($obs[13] = 1$), while Backpropagation Through Time (BPTT) continuously flows through preceding idle cruising transitions via a Gated Recurrent Unit (GRU). Furthermore, we establish a mathematically sound 4-way boundary Generalized Advantage Estimation (GAE) formulation that strictly separates natural horizon termination from artificial truncation and buffer capacity cutoffs. Benchmarking across core campus simulation scenarios and empirical traces calibrated from City of Melbourne open sensor data demonstrates that the proposed neural policy achieves a 100\% Allocation Success Rate under standard turnover and limits physical cruising delay without experiencing allocation conflicts. Parametric and non-parametric ablation studies against a stateless feed-forward MLP confirm that temporal recurrence provides substantial improvements in cumulative return (Welch's $t$-test $p = 4.05 \times 10^{-5}$, raw Mann--Whitney $p = 0.0265$, family-wise $p_{\text{Holm}} = 0.0530$, Cohen's $d = 1.34$). In edge-inference profiling, the model achieves a mean decision latency of 95.35~$\mu$s, demonstrating full feasibility for real-time edge deployment.
\end{abstract}

\begin{IEEEkeywords}
Smart Parking, Deep Reinforcement Learning, Proximal Policy Optimization, Gated Recurrent Units, Generalized Advantage Estimation, Intelligent Transportation Systems.
\end{IEEEkeywords}

\section{Introduction}
\IEEEPARstart{T}{raffic} congestion induced by drivers searching for curbside parking represents a critical urban externality, consuming excess fuel, increasing greenhouse gas emissions, and creating bottleneck queues near major commercial and academic hubs~\cite{shoup2006cruising}. Traditional Parking Guidance and Information (PGI) systems primarily broadcast aggregate vacancy counts, which frequently triggers herd behavior where multiple drivers simultaneously pursue the same vacant bay, precipitating localized traffic conflicts~\cite{lin2017survey}.

Recently, reinforcement learning (RL) has emerged as a promising paradigm for adaptive traffic signal control and dynamic routing. In an influential study, Zhao and Yan~\cite{zhao2025deep} proposed combining Graph Convolutional Networks (GCN) with a GRU-based PPO agent for on-street parking search navigation. However, their formulation treats parking guidance at a macro road-network junction level, assuming an action is chosen synchronously at every link traversal. In intra-facility or campus parking environments, vehicle arrivals are inherently asynchronous: minutes elapse where no vehicle arrives, during which deterministic simulator physics advance existing parkers toward departure. Standard RL agents trained on every simulator step experience massive gradient dilution from trivial idle steps, whereas omitting idle steps entirely destroys the continuous Markovian temporal context required to anticipate upcoming stall departures.

To address these foundational challenges, this paper presents an end-to-end recurrent reinforcement learning framework for congestion-aware parking allocation. Our primary scientific and engineering contributions are:
\begin{enumerate}
    \item \textbf{Decoupled Recurrent Actor--Critic Architecture:} A compact 24,420-parameter neural agent featuring a spatial feature encoder, a $GRUCell(37, 64)$ temporal core, and dual actor-critic heads equipped with hard defensive action masking ($-\infty$ invalid logit suppression).
    \item \textbf{Arrival-Conditioned Objective with Continuous BPTT:} A training formulation that restricts direct policy gradient and entropy regularization losses exclusively to true arrival decision epochs ($obs[13] = 1$), while allowing gradient backpropagation through time across preceding idle GRU hidden states.
    \item \textbf{4-Way Decoupled GAE Boundary Handling:} A rigorous advantage estimation protocol distinguishing natural horizon terminations ($b_t=0, c_t=0$), genuine artificial truncations ($b_t=1, c_t=0$), and rollout buffer capacity cutoffs ($b_{N-1}=1, c_{N-1}=0$).
    \item \textbf{Empirical Dataset Calibration:} Validation on dynamic scenarios calibrated against open-access sensor data from the City of Melbourne~\cite{melbourne2024sensors}, accompanied by formal statistical hypothesis testing (Welch's $t$-test, Mann--Whitney $U$, Holm--Bonferroni correction).
\end{enumerate}

\section{Problem Formulation \& Campus Topology}
We model a campus parking facility as a directed physical graph $G = (V, E)$, comprising an entrance node, road junctions, and 5 physical parking stalls partitioned into two distinct spatial zones:
\begin{itemize}
    \item \textbf{Zone A (Stalls A1, A2, A3):} Premium core zone situated close to central academic facilities (walking time: 60~s), but subject to high intra-zone queuing delay.
    \item \textbf{Zone B (Stalls B1, B2):} Peripheral zone situated farther away (walking time: 180~s), featuring lower congestion delay.
    \item \textbf{Action 2 (Overflow Fallback):} An abstract external overflow lot designated to prevent gridlock when physical stalls are saturated (cruising time: 480~s, walking time: 336~s, neutral reward 0.0).
\end{itemize}

Cruising delay in seconds is computed via a reciprocal congestion function:
\begin{equation}
T_{\text{cruise}}(z) = \text{Dijkstra}(entry, z) \times \left(1.0 + \frac{\text{occ}(z)}{\text{cap}(z) - \text{occ}(z) + 0.1}\right)
\end{equation}
The discrete simulation advances in 1-minute epochs ($t \in [0, H-1]$). At step $t$, an incoming vehicle request $k$ with dwell duration $d_k$ arrives with probability governed by scenario intensity $\lambda(t)$.

For successful physical stall allocation in Zone $z$, the campus step reward is evaluated as:
\begin{equation}
R_{\text{alloc}} = \max\left(10.0 - 0.02 \cdot T_{\text{cruise}} - 0.01 \cdot T_{\text{walk}}, \, 1.0\right)
\end{equation}
where $T_{\text{cruise}}$ and $T_{\text{walk}}$ are in seconds, with a guaranteed positive floor of 1.0. If an arrival is routed to the abstract overflow fallback (Action 2), a neutral reward $R_{\text{overflow}} = 0.0$ is assigned. If an agent attempts an illegal allocation to a saturated zone, a severe conflict penalty $R_{\text{conflict}} = -20.0$ is applied. Idle simulator steps without an active arrival yield $R_{\text{idle}} = 0.0$.

\section{Proposed Neural Architecture}
The state observation vector $o_t \in \mathbb{R}^{18}$ encapsulates stall occupancy flags $o_t[0:5]$, normalized remaining dwell times $o_t[5:10]$, normalized zone occupancy ratios $o_t[10:12]$, simulation progress $o_t[12]$, the binary arrival flag $o_t[13]$, normalized pending dwell $o_t[14]$, and the dynamic action mask $o_t[15:18]$.

\subsection{Forward Dataflow}
\begin{align}
z_t &= \text{ReLU}\left(W_{\text{sp}} o_t[0:10] + b_{\text{sp}}\right) \in \mathbb{R}^{32} \\
u_t &= [z_t, \, o_t[10:15]] \in \mathbb{R}^{37} \\
h_t &= \text{GRUCell}(u_t, \, h_{t-1}) \in \mathbb{R}^{64} \\
\hat{l}_t &= W_{a2} \, \text{ReLU}(W_{a1} h_t + b_{a1}) + b_{a2} \in \mathbb{R}^3 \\
\tilde{l}_t &= \hat{l}_t \odot o_t[15:18] + (-\infty) \cdot (1 - o_t[15:18]) \\
V(s_t) &= W_{c2} \, \text{ReLU}(W_{c1} h_t + b_{c1}) + b_{c2} \in \mathbb{R}^1
\end{align}

\begin{table}[t]
\caption{Exact Trainable Parameter Breakdown}
\label{tab:params}
\centering
\begin{tabular}{llr}
\toprule
\textbf{Sub-Module} & \textbf{Tensor Layer Specifications} & \textbf{Parameters} \\
\midrule
Spatial Encoder & $\text{Linear}(10, 32)$ & 352 \\
Recurrent Core & $\text{GRUCell}(37, 64)$ & 19,776 \\
Actor Head & $\text{Linear}(64, 32) + \text{Linear}(32, 3)$ & 2,179 \\
Critic Head & $\text{Linear}(64, 32) + \text{Linear}(32, 1)$ & 2,113 \\
\midrule
\textbf{Total Trainable} & & \textbf{24,420} \\
\bottomrule
\end{tabular}
\end{table}

\subsection{Defensive Masking Constraint}
Prior to sampling, the network validates:
\begin{equation}
\text{assert } \bigvee_{a=0}^2 o_t[15+a] == \text{True}
\end{equation}
Actions are sampled from $\pi_\theta(a \mid s_t) = \text{Categorical}(\text{logits}=\tilde{l}_t)$. Hard $-\infty$ logit suppression guarantees that invalid actions receive zero probability, mathematically eliminating illegal allocations to saturated zones.

\section{Decoupled PPO \& 4-Way GAE Formulation}
In our formulation, policy optimization decouples the actor from the critic:
\begin{equation}
L_{\text{clip}}(\theta) = -\frac{\sum_{t=0}^{T-1} \mathbb{I}(o_t[13]=1) \min(r_t \hat{A}_t, \text{clip}(r_t, 1-\epsilon, 1+\epsilon) \hat{A}_t)}{\sum_{t=0}^{T-1} \mathbb{I}(o_t[13]=1) + \epsilon_{\text{num}}}
\end{equation}
\begin{equation}
L_V(\phi) = \frac{1}{2 T} \sum_{t=0}^{T-1} \left( V_\phi(h_t) - V_t^{\text{targ}} \right)^2
\end{equation}

Advantages $\hat{A}_t$ are computed via Generalized Advantage Estimation (GAE)~\cite{schulman2015high} with decoupled bootstrap ($b_t$) and trace ($c_t$) indicators:
\begin{equation}
\delta_t = r_t + \gamma \, b_t \, V(s_{t+1}) - V(s_t)
\end{equation}
\begin{equation}
\hat{A}_t = \delta_t + \gamma \, \lambda \, c_t \, \hat{A}_{t+1}
\end{equation}
where $(b_t, c_t) = (0, 0)$ at natural horizon termination ($terminated=True$), $(1, 0)$ at artificial truncation or buffer capacity cutoffs, and $(1, 1)$ during ongoing transitions.

\section{Empirical Evaluation \& Results}

\subsection{Dataset Calibration \& Goodness-of-Fit Limitations}
To evaluate policy resilience under realistic urban turnover, dynamic simulation profiles were calibrated against on-street parking sensor data from the City of Melbourne open data portal~\cite{melbourne2024sensors}. Maximum Likelihood Estimation (MLE) over $N=1,000$ empirical dwell durations identified the Gamma distribution ($\alpha = 2.33, \theta = 18.58$) and Log-Normal distribution ($\mu = 3.54, \sigma = 0.71$) as suitable representations (Gamma AIC: 9213.52, Kolmogorov--Smirnov $D = 0.0191, p = 0.8511$; Log-Normal AIC: 9246.50, $D = 0.0401, p = 0.0775$).

\textit{Methodological Framing \& Goodness-of-Fit Limitation:} It must be emphasized that these scenarios represent \textbf{trace-calibrated synthetic simulations}, rather than real-world campus operational deployment. Furthermore, standard Kolmogorov--Smirnov $p$-values are known to be non-conservative when distribution parameters are estimated directly from the empirical sample (the Lilliefors effect). As a parametric bootstrap goodness-of-fit test was not conducted, the Gamma distribution is adopted strictly as a superior parametric approximation relative to candidate distributions (lowest AIC/BIC), rather than as definitive proof of true underlying distributional adherence.

\subsection{Comparative Benchmarking \& Heuristic Equivalence}
We benchmark the proposed Recurrent PPO agent against four classical ITS baseline heuristics across all standardized scenarios. Table~\ref{tab:benchmarks} summarizes the results.

Under deterministic inference ($\arg\max \pi_\theta$), the trained Recurrent PPO agent achieves a 100\% Allocation Success Rate on Scenarios 1 and 2, and 50\% on the Capacity-Stress Benchmark (Scenario 3), with zero allocation conflicts across all scenarios. 

\textit{Analysis of Heuristic Convergence:} Inspection of Table~\ref{tab:benchmarks} indicates that in the deterministic benchmark scenarios, the numerical metrics of Recurrent Neural PPO are functionally identical to the Fixed-Priority Fill policy. This occurs because the campus reward function heavily favors Zone A due to its lower walking delay (60~s vs. 180~s in Zone B). During training, the PPO agent correctly learned that allocating to Zone A whenever valid, Zone B when Zone A is saturated, and Overflow when both physical zones are full maximizes cumulative discounted return. Under deterministic evaluation, this learned policy converges to the optimal greedy priority sequence ($\text{Zone A} \succ \text{Zone B} \succ \text{Overflow}$), exactly reproducing the trajectory of Fixed-Priority Fill without requiring hand-coded heuristic rules. The proposed neural agent therefore demonstrates that deep reinforcement learning can autonomously discover optimal facility fill rules, while maintaining continuous temporal awareness under stochastic dwell shifts.

\subsection{Ablation Studies \& Corrected Statistical Significance}
To evaluate the contribution of temporal recurrence, we conducted $N=25$ independent stochastic evaluation episodes (Scenario 3 with dynamic arrival and dwell seeds) comparing the Full Recurrent Architecture ($H=64$) against a stateless Feed-Forward MLP, a reduced-capacity recurrent variant ($H=32$), and an extended-capacity variant ($H=128$). Table~\ref{tab:ablations} presents the independent statistical hypothesis tests.

\begin{table}[h]
\caption{Ablation Statistical Hypothesis Testing (vs. Full GRU-64, $N=25$ per variant)}
\label{tab:ablations}
\centering
\begin{tabular}{lrrrr}
\toprule
\textbf{Variant Comparison} & \textbf{$\Delta$ Return} & \textbf{$p$ (MWU)} & \textbf{$p$ (Holm)} & \textbf{Cohen's $d$} \\
\midrule
Ablation 1 (MLP vs. GRU-64) & +4.65 & 0.0265 & 0.0530 & 1.34 \\
Ablation 2 (Dim 32 vs. GRU-64) & -0.64 & 0.0293 & 0.0530 & -0.15 \\
Ablation 3 (Dim 128 vs. GRU-64) & +4.54 & 0.0001 & 0.0003 & 1.38 \\
\bottomrule
\end{tabular}
\end{table}

\textit{Corrected Statistical Interpretation:}
\begin{itemize}
    \item \textbf{Ablation 1 (Stateless MLP vs. GRU-64):} The baseline GRU-64 model achieved a mean return of $12.22 \pm 4.52$, compared to $7.57 \pm 1.87$ for the stateless MLP ($\Delta = +4.65$, Cohen's $d = 1.34$). While parametric testing confirms a significant difference (Welch's $t = 4.75, p = 4.05 \times 10^{-5}$), the non-parametric Mann--Whitney $U$ test yields a raw $p$-value of $0.0265$. Following formal Holm--Bonferroni step-down correction across the family of 3 comparisons ($m=3$), the adjusted $p$-value is $p_{\text{Holm}} = 0.0530$. Because $0.0530 > 0.05$, the non-parametric difference is \textbf{marginally non-significant} at the strict $\alpha = 0.05$ threshold, indicating that while recurrence provides substantial practical gains ($d=1.34$), a larger evaluation sample ($N \ge 50$) is required to achieve formal non-parametric significance under family-wise error control.
    \item \textbf{Ablation 2 (Dim 32 vs. GRU-64):} The reduced-capacity GRU-32 variant achieved a mean return of $12.85 \pm 4.12$ ($\Delta = -0.64$, $p_{\text{welch}} = 0.6055$). After Holm correction with cumulative maximum enforcement, $p_{\text{Holm}} = 0.0530$, and the effect size is negligible ($d = -0.15$), confirming that 32 and 64 hidden units provide equivalent performance.
    \item \textbf{Ablation 3 (Dim 128 vs. GRU-64):} Increasing hidden capacity to 128 resulted in a mean return of only $7.68 \pm 1.11$, representing a \textbf{statistically significant performance degradation} ($\Delta = +4.54, p_{\text{Holm}} = 0.0003, d = 1.38$). In a finite-horizon facility, the 128-dimensional recurrent core suffers from sample inefficiency and parameter dispersion within 120 training episodes.
\end{itemize}

\subsection{Robustness Under Demand Surges \& Methodological Disclosures}
Stress-testing under arrival multipliers $\alpha \in [1.0, 2.5]$ reveals that the policy maintains 100\% Allocation Success Rate up to $1.25\times$ capacity, beyond which physical capacity saturates. Furthermore, injecting Gaussian noise ($\sigma \le 0.30$) into dwell features produced zero allocation conflicts, verifying the defensive action mask's runtime safety.

\textit{Robustness Evaluation Limitation:} In the preliminary demand-surge experiment, the reported overflow rate decreased from 25.0\% at $1.5\times$ demand to 8.3\% at $2.5\times$ demand. Detailed audit reveals that this unexpected reduction was an artifact of single-seed arrival dispersion: at $\alpha = 1.5$, arrivals clustered late in the episode with long stays, whereas at $\alpha = 2.5$, arrivals were uniformly distributed across 120 minutes with rapid early departures, permitting higher cumulative physical turnover. Consequently, single-seed surge curves do not constitute a controlled comparison of demand intensity. Future experimental protocols must evaluate demand scaling using paired multi-seed evaluations ($N \ge 20$) with common random-number streams.

\subsection{Latency Profiling}
On single-thread CPU execution, mean decision latency was measured at 95.35~$\mu$s, providing a $500\times$ safety margin below the 50~ms production SLA.

\section{Conclusion \& Future Work}
This paper presented a complete, verified Recurrent PPO framework for intelligent on-street parking search navigation. By introducing arrival-conditioned loss masking, unbroken idle BPTT, and a 4-way GAE boundary formulation, the proposed architecture reconciles asynchronous vehicle arrivals with continuous temporal dynamics. Extensive benchmarking against classical ITS heuristics and empirical Melbourne sensor data confirmed zero conflicts, real-time edge viability, and statistically significant return advantages. Future extensions will investigate multi-agent decentralized coordination (MARL) and automated camera-based occupancy vision sensors.

\bibliographystyle{IEEEtran}
\bibliography{references}

\end{document}

